\documentclass[10pt,a4paper]{amsart}
\usepackage[margin=19mm]{geometry}
\usepackage{amsmath,amssymb,mathtools}
\usepackage[scr=rsfs,cal=euler]{mathalfa}
\usepackage{xfrac}
\usepackage[unicode,hidelinks,bookmarks=false]{hyperref}
\newcommand{\ie}{i.e.\ }
\newcommand{\cf}{cf.\ }
\newcommand{\resp}{respectively\ }
\newcommand{\Subsection}{\subsection}
\providecommand{\nobreakdash}{\nobreak\hbox{-}}
\setlength{\emergencystretch}{2em}
\multlinegap0pt
\usepackage{xfrac}
\usepackage{mathtools}
\usepackage{nicefrac}
\usepackage{csquotes}
\usepackage{xcolor}
\usepackage{bm}
\usepackage{caption}
\usepackage{amssymb}
\usepackage[shortlabels]{enumitem}
\usepackage{dsfont}
\usepackage{tikz}
\usepackage{booktabs}
\usepackage{algorithm}\usepackage{algpseudocode}
\newtheorem{assumption}{Assumption}
\newtheorem{lemma}{Lemma}
\newtheorem{theorem}{Theorem}
\usepackage{cleveref}
\crefname{assumption}{Assumption}{Assumptions}
\crefname{lemma}{Lemma}{Lemmas}
\crefname{theorem}{Theorem}{Theorems}
\newcommand{\N}{\mathbb N}
\def\R{\mathbb{R}}
\newcommand{\Res}{\calR}%{\text{R}}
\newcommand{\bu}{\mathbf{u}}
\newcommand{\calL}{\mathcal{L}}
\newcommand{\calR}{\mathcal{R}}
\newcommand{\ddt}{\partial_t}
\newcommand{\dt}{\mathrm{d}t}
\newcommand{\linspan}{\mathop{\rm span}\nolimits}
\newcommand{\mU}{{\mathscr U}}
\newcommand{\mY}{{\mathscr Y}}
\newcommand{\norm}[2]{\left|#1\right|_{#2}}
\newcommand{\reduce}[1]{{{#1}^r}}
\newcommand{\scalprod}[3]{\langle#1, #2\rangle_{#3}}
\newcommand{\tildeyint}{\tilde y_{\mathrm{in}}}
\newcommand{\tint}{t_{\mathrm{in}}}
\newcommand{\yint}{y_{\mathrm{in}}}
\title{Stabilization of parabolic time-varying PDEs using certified reduced-order receding horizon control}
\author{Behzad Azmi, Michael Kartmann, Stefan Volkwein}
\date{Source: arXiv:2508.16801v2 (CC BY 4.0)}
\begin{document}
\maketitle

\noindent\emph{Source and licence.} Stabilization of parabolic time-varying PDEs using certified reduced-order receding horizon control. Version arXiv:2508.16801v2 (CC BY 4.0), distributed by arXiv under the Creative Commons Attribution 4.0 International licence, http://creativecommons.org/licenses/by/4.0/, as recorded in the arXiv licence metadata for this exact version and shown on the abstract page https://arxiv.org/abs/2508.16801v2. Full licence notice: \texttt{licenses/arxiv-2508.16801-CC-BY-4.0.txt}.\par\smallskip

\noindent\textit{Editorial scope.} Counted unit: the article's theorem theo:APP\_convergence (source0.tex lines 2304--2310). The article's complete original proof of it is delivered with it (source0.tex lines 2312--2399, one \texttt{proof} environment of 88 lines). No mathematics is written, repaired or re-derived: the statement and the proof are the article's own lines, in the article's own notation, and no retained line is substituted or re-flowed.  Retained uncounted as the source-local context the proof needs: lines 633--638; lines 642--651; lines 647--649; lines 1002--1006; lines 1045--1057; lines 1080--1084; lines 1121--1123; lines 1127--1152; lines 1204--1231; lines 1850--1854; lines 2233--2253; lines 2239--2241; lines 2246--2248.  Nothing was omitted from any retained span, so the package carries no omission record.  No retained line contains a citation, so the wrapper carries no bibliography and no bibliography entry.  No retained line contains a figure environment, an image-inclusion command or drawing code, so no image is delivered and the package declares no asset dependency.  Editorial formatting only, in addition: an AMS \texttt{amsart} preamble that restates the article's own declarations and macros used by the retained lines, namely the environments \texttt{assumption}, \texttt{lemma}, \texttt{theorem} on separate counters, together with the standard packages those lines need.  The retained environments print in this document as Assumption 1, Lemma 1, Theorem 1 in order of appearance, whereas the article prints the same items with the numbering of its own full document.  The complete native source of the article as distributed in the arXiv e-print for this exact version is bundled unchanged as \texttt{sources/183723/source0.tex}.
\begin{align}
    \label{eq:LTV}
    %
    \tag{$\mathrm{FOM}_T(\tint,\yint)$}
    \ddt y(t)+A(t)y(t)=B(t)\bu(t), \quad t \in  ( \tint,\tint+T),\quad y(\tint)=\yint.
\end{align}

\begin{assumption}\label{ass:pde}
    %
    We assume that $A\in L^\infty(0,\infty;\mathcal{L}(V,V'))$, $B \in  L^\infty(0,\infty;\mathcal{L}(U,V'))$
     and the existence of constants $\eta_V>0$, $\eta_H\geq 0$ such that
    %
    \begin{align}\label{eq:weak_coercivity}
        \scalprod{A(t)v}{v}{V',V} \geq \eta_V\norm{v}{V}^2 - \eta_H\norm{v}{H}^2 \quad \text{for all } v\in V,\  t\geq 0.
    \end{align}
    %
\end{assumption}

    \begin{align}
        \scalprod{A(t)v}{v}{V',V} \geq \eta_V\norm{v}{V}^2 - \eta_H\norm{v}{H}^2 \quad \text{for all } v\in V,\  t\geq 0.
    \end{align}

    \begin{align}
        \label{eq:FOM_opsys}
            -\ddt p(t)+A'(t)p(t)= y(t), \quad t \in  (\tint,\tint+T),\quad 
            p(\tint+T)=0.
    \end{align}

\begin{equation}
    \label{eq:reducedLTV}
    \tag{$\text{ROM}^r_{T}(\tint, \tildeyint)$}
     \ddt {\reduce {y}}(t)+A(t){{\reduce {y}}}(t)= {B}(t)\bu(t)\ \text{in } V_r', \ t \in  ( \tint,\tint+T),\ 
           \reduce{y}(\tint)
           =\Pi^H_{V_r} \tildeyint,
    % \left\{
    % \begin{array}{rll}
    %     \ddt \reduce{y}(t)+{A}(t)\reduce{y}(t)&={B}(t)\bu(t)&\text{in } V_r'\text{ for } t \in (\tint,\tint+T),\\
    %     \reduce{y}(\tint)&=\Pi^H_{V_r} \tildeyint,&
    % \end{array}
    % \right.
\end{equation}

    \begin{align}
        \label{eq:ROM_opsys}
            -\ddt {\reduce {p}}(t)+A'(t){\reduce {p}}(t)= { \tilde{y}}(t)\ \text{in } V_r', \quad t \in  ( \tint,\tint+T),\quad 
            {\reduce {p}}(\tint+T)=0,
    \end{align}

\begin{equation}\label{eq:perturb_init}
    \norm{ \yint-\tildeyint}{H} \leq \Delta_{\yint} \quad \text{for } \Delta_{\yint}\geq 0.
\end{equation}

\begin{lemma}[State a posteriori estimator]
    \label{Lemma:ResBasStateError}
    Let \cref{ass:pde} be valid and
    let $y=y(\bu,\tint,\yint)\in\mY_T(\tint)$ and $y^r=y^r(\bu^r ,\tint,\tildeyint)\in\mY^r_T(\tint)$ be the solution of \eqref{eq:LTV} and \eqref{eq:reducedLTV} for $\yint,\tildeyint \in H$ and $\bu,\bu^r\in\mU_T(\tint)$, respectively. Assume the data satisfies \eqref{eq:perturb_init} and ${\norm{\bu-\reduce\bu}{\mU_T(\tint)}}\leq\Delta_\bu$
    for $\Delta_\bu\geq 0$. Define the state error and residual as
    %
    \begin{equation}\nonumber
        e_y\coloneqq y-\reduce y, \quad \Res_y(y^r,\bu^r)(t)\coloneqq B(t)\bu^r(t)-A(t)y^r(t)-\partial_t y^r(t) \in V' \text{ for }t\in(0,T).
    \end{equation}
    %
    Then, we have the \emph{a posteriori} error bound for $t\in [\tint, \tint+T]$
    %
    \begin{align}\label{eq:state_apost_CHL2V}
        &\norm{e_y(t)}{H}^2+\norm{e_y}{L^2(\tint,t;V)}^2\leq \Delta_y^2(t;\Delta_\bu, \Delta_{\yint}, y^r, \bu^r)
    \end{align}
    %
    with
    %
    \begin{align*}
        \Delta^2_{y}(t)&\coloneqq C_{1,y}(t)\Big(\norm{e^{-\eta_H (\cdot-\tint) }B}{L^\infty(\tint,t;\calL(U,V'))}^2{\Delta_\bu^2} + \norm{e^{-\eta_H (\cdot-\tint) }\Res_y(\reduce y, \reduce \bu)}{L^2(\tint,t;V')}^2\Big)\\
        &\hspace{1mm}\quad+ C_{2,y}(t){\big( \Delta_{\yint} + \norm{\tildeyint - \Pi^H_{V_r}\tildeyint }{H} \big)^2}
    \end{align*}
    %
    for the constants $C_{1,y}(t)\coloneqq \nicefrac{2e^{2\eta_H(t-\tint)}}{\min\{1,\eta_V\}\eta_V}$, $C_{2,y}(t)\coloneqq \nicefrac{e^{2\eta_H(t-\tint)}}{\min\{1,\eta_V\}}$. If $t=\tint+T$, we simply write $\Delta_y^2(\Delta_\bu, \Delta_{\yint}, y^r, \bu^r)$ instead of $\Delta_y^2(\tint+T;\Delta_\bu, \Delta_{\yint}, y^r, \bu^r)$.
    %
\end{lemma}

\begin{lemma}[Adjoint state a posteriori estimator]
    \label{lem:ResAdjest}
    Suppose that \cref{ass:pde} holds. Let $p=p(y)\in\mY_T(\tint)$ and $p^r=p^r(\tilde y)\in\mY^r_T(\tint)$ be the solutions of \eqref{eq:FOM_opsys} and \eqref{eq:ROM_opsys} for $y,\tilde y\in L^2(\tint,\tint+T;H)$, respectively. Assume that $\norm{y-\tilde y}{L^2(\tint,\tint+T;H)} \leq \Delta_y $
    for $\Delta_y\geq 0$. Define the adjoint state error and residual as
    %
    \begin{equation*}
        e_p\coloneqq p-p^r, \quad \Res_p(\tilde y, p^r)(t)\coloneqq\tilde y(t)-A'(t)p^r(t)+\dot p^r(t)\in V'.
    \end{equation*}
    %
    Then, we have the \emph{a posteriori} error bound
    %
    \begin{align}\label{eq:adjstate_apost_CHL2V}
        &\norm{e_p(\tint)}{H}^2+\norm{e_p}{L^2(\tint,\tint+T;V)}^2\leq \Delta_p^2(\Delta_y,\tilde y, p^r)
    \end{align}
    %    
    with
    %
    \begin{align}\label{eq:adjstate_apost_CHL2V_struc}
        \Delta_{p}^2{(\Delta_y,\tilde y, p^r)}  \coloneqq & C_{p}\Big(\Delta_y^2 + \norm{e^{\eta_H (\cdot-\tint-T) }\Res_p(\tilde y, \reduce p)}{L^2(\tint,\tint+T;V')}^2\Big)
    \end{align}
    %
    and $C_{p}\coloneqq\nicefrac{2e^{2\eta_H T}}{\min(1,\eta_V)\eta_V}$. Further, we have the improved estimate
    %
    \begin{align}\label{eq:adjstate_apost_Ht}
        &\norm{e_p(\tint)}{H}^2\leq \tfrac{1}{2} \Delta_{p}^2 {(\Delta_y,\tilde y, p^r)}.%\Delta^2_{p,H}( \tilde y, p^r)\coloneqq \tfrac{1}{2} \Delta_{p}^2 (\tilde y,p^r).
    \end{align}
    %
\end{lemma}

\begin{align}
    \label{eqn:PODminimization}
     \min_{V_r=\linspan(v_i)_{i=1}^r\subseteq V} \sum_{s\in S}\norm{s-\Pi^V_{V_r}s }{L^2(t_k,t_k+T;V)}^2\mathrm dt
     \ \  \text{s.t.}\ \ \langle v_i, v_j\rangle_V = \delta_{ij} \ (i,j = 1,\ldots, r),
\end{align}

\begin{lemma}[Error estimator equivalence in $\mY_T(\tint)$]\label{lemma:app_error_equivalence}
    %
    Let \cref{ass:pde} hold and let $\tint\in\R_{\geq0}$, $T\in \R_{>0}$.
    %
    \begin{enumerate}
        \item In the situation of \cref{Lemma:ResBasStateError}, let $y=y(\bu,\yint)\in \mY_T(\tint),\  y^r=y^r(\bu,\yint)\in \mY^r_T(\tint)$ be the solution of \eqref{eq:LTV} and \eqref{eq:reducedLTV}, respectively, for $\tildeyint=\yint\in H$, $\bu\in \mU_T(\tint)$. For $e_y = y-y^r$ it holds
     \begin{equation}\label{eq:state_errorest_equiv}
     c \Delta^2_{y}(0,0, y^r,\bu)\leq \norm{e_y(\tint+T)}{H}^2+\norm{e_y}{\mY_T(\tint)}^2\leq C\Delta^2_y(0,0,y^r, \bu),
    \end{equation}
    %
    for constants $C,c>0$ independent of $(y,y^r,\bu, \yint,r)$.
    \item In the situation of \cref{lem:ResAdjest}, let $p=p(\tilde y)\in  \mY_T(\tint)$ and $p^r=p^r(\tilde y)\in \mY^r_T(\tint)$ be the solution of the FOM \eqref{eq:FOM_opsys} and ROM \eqref{eq:ROM_opsys} adjoint equation, respectively, for $\tilde y\in L^2(\tint,\tint+T;H)$. For $e_p = p-p^r$ it holds
    %
    \begin{equation}\label{eq:adstate_errorest_equiv}
    \tilde c\Delta^2_{p}(0, p^r, \tilde y) \leq \norm{e_p(\tint)}{H}^2+\norm{e_p}{\mY_T(\tint)}^2\leq \tilde C\Delta^2_{p}(0, p^r, \tilde y),
    \end{equation}
    %
    for constants $\tilde C,\tilde c>0$ independent of $(p,p^r,\tilde y,r)$.
    \end{enumerate}
    %
\end{lemma}

     \begin{equation}
     c \Delta^2_{y}(0,0, y^r,\bu)\leq \norm{e_y(\tint+T)}{H}^2+\norm{e_y}{\mY_T(\tint)}^2\leq C\Delta^2_y(0,0,y^r, \bu),
    \end{equation}

    \begin{equation}
    \tilde c\Delta^2_{p}(0, p^r, \tilde y) \leq \norm{e_p(\tint)}{H}^2+\norm{e_p}{\mY_T(\tint)}^2\leq \tilde C\Delta^2_{p}(0, p^r, \tilde y),
    \end{equation}

\begin{theorem}[$\mY_T(\tint)$-convergence of the state and adjont state]\label{theo:APP_convergence}
%
In the situation of \ \cref{lemma:app_error_equivalence}, 
consider an orthonormal basis $(v_n)_{n\in \N}$ of $V$. We set $V_r\coloneqq \text{span}\{v_1,\ldots,v_r\}$. Then we have $\norm{ e_y}{\mY_T(\tint)}\to 0$, $\norm{ e_p}{\mY_T(\tint)}\to 0$ as $r\to \infty$. 
%,$\Delta_y^2(0,0,y^r,\bu)\to 0$ and $\Delta_p(0, p^r, \tilde y)\to 0$ as $r\to \infty$.
%
\end{theorem}

\begin{proof}
    %
    We show $\norm{e_y}{\mY_T(\tint)}\to 0$, as the claim for $e_p$ follows by similar arguments. 
    % From there, the decay of the error estimators follows from the equivalences \eqref{eq:state_errorest_equiv}, \eqref{eq:adstate_errorest_equiv} in \cref{lemma:app_error_equivalence} and the embedding $\mY_T(\tint)\hookrightarrow C([\tint,\tint+T];H)$.
    Consider the orthogonal projection operator $\Pi^V_{V_r}:V\to V_r$ defined below \eqref{eqn:PODminimization}. We decompose the error according to $ y^r-y = y^r-\Pi^V_{V_r} y+\Pi^V_{V_r} y-y\eqqcolon e_2 + e_1$. Consider $e_1$ first. It holds $\ddt(\Pi^V_{V_r} y)= \dot y \circ \Pi^V_{V_r} $ (since $\Pi^V_{V_r}$ is self-adjoint as an orthogonal projection) and therefore
    %
    \begin{align*}
        \norm{e_1}{\mY(\tint)}^2=&\norm{\Pi^V_{V_r} y-y}{L^2(V)}^2+\norm{\dot y \circ \Pi^V_{V_r} -\dot y}{L^2(V')}^2\\
        % = &\|\Pi^V_{V_r} y-y\|_{L^2(V)}^2+\sup_{v\in L^2(V), \|v\|_{L^2(V)}=1}\int_{\tint}   ^{\tint+T}|\dot y(t) \circ \Pi^V_{V_r}(v) -\dot y(t)(v)|\ \dt\\
         %= &\norm{\Pi^V_{V_r} y-y}{L^2(V)}^2+\sup_{v\in L^2(V), \norm{v}{L^2(V)}=1}|\langle \dot y\circ \Pi^V_{V_r} -\dot y, v\rangle_{L^2(V'),L^2(V)}|^2\\
        %=& \norm{\Pi^V_{V_r} y-y}{L^2(V)}^2+\sup_{v\in L^2(V), \norm{v}{L^2(V)}=1}|\langle \dot y,  \Pi^V_{V_r}v-v\rangle_{L^2(V'),L^2(V)}|^2\\
        \leq &  \norm{\Pi^V_{V_r} y-y}{L^2(V)}^2+\norm{\dot y}{L^2(V')}^2\sup_{v\in L^2(V), \norm{v}{L^2(V)}=1}\norm{\Pi^V_{V_r}v-v}{L^2(V)}^2\to 0 \ (r\to \infty).
    \end{align*}
    %
    Now we turn to $e_2$. It holds for $v\in V_r$ and almost all $t\in (\tint,\tint+T)$
    %
    \begin{align}\label{eq:theo:galerkin_conv_help0}
        \langle\dot e_2(t),v\rangle_{V',V}+\langle A(t)e_2(t), v\rangle_{V',V}=& \langle-\dot e_1(t),v\rangle_{V',V}+\langle -A(t)e_1(t), v\rangle_{V',V}\\ \nonumber
        \leq & \norm{\dot e_1(t)}{V'}\norm{v}{V}+\norm{A}{L^\infty}\norm{e_1(t)}{V}\norm{v}{V}
    \end{align}
    %
    By choosing $v=e_2(t)$ and using the weak coercivity of the operator $A$ as stated in \eqref{eq:weak_coercivity}, along with Young's inequality, we arrive at
    %
    \begin{align}\label{eq:theo:galerkin_conv_help1}
         \tfrac{1}{2}\norm{\dot e_2(t)}{H}^2+\tfrac{\eta_V}{2}\norm{e_2(t)}{V}^2
        \leq & \tfrac{1}{\eta_V}\norm{\dot e_1(t)}{V'}^2+\tfrac{\norm{A}{L^\infty}}{\eta_V} \norm{e_1(t)}{V}^2+\eta_H\norm{e_2(t)}{H}^2
    \end{align}
    %
    %
    By applying Gronwall's Lemma and integration over $(\tint,\tint+T)$ it follows that
    %
    % \begin{equation}
    %     \nonumber 
    %     \norm{e_2(t)}{H}^2\leq c(T)\bigg[  \norm{e_2(\tint)}{H}^2+\norm{\dot e_1(t)}{V'}^2+\norm{e_1(t)}{V}^2\bigg].
    %     % \norm{e_2(t)}{H}^2\leq e^{2\eta_HT}\bigg[  \norm{e_2(\tint)}{H}^2+\tfrac{1}{\eta_V}\norm{\dot e_1(t)}{V'}^2+\tfrac{\norm{A}{L^\infty}}{\eta_V} \norm{e_1(t)}{V}^2\bigg].
    % \end{equation}
    % %
    % Plugging this in \eqref{eq:theo:galerkin_conv_help1}, we obtain
    % %
    % \begin{align*}
    %     \norm{e_2(t)}{V}^2\leq c\norm{\dot e_1(t)}{V'}^2+c \norm{e_1(t)}{V}^2+c\norm{e_2(\tint)}{H}^2.
    % \end{align*}
    % %
    % Integration over $(\tint,\tint+T)$ gives 
    % %
    \begin{align}\label{eq:theo:galerkin_conv_help2}
        \norm{e_2}{L^2(V)}^2\leq c\norm{\dot e_1}{\mY_T(\tint)}^2+c\norm{e_2(\tint)}{H}^2.
    \end{align}
    %
    for a generic constant $c=c(T)$.
    For $e_2(\tint)$, we obtain, using $\mY_T(\tint)\hookrightarrow C([\tint,\tint+T];H)$, that
    %
    \begin{align*}
      \norm{e_2(\tint)}{H}^2=&\norm{\big(y^r-y+y-\Pi^V_{V_r} y\big)(\tint)}{H}^2\leq  \norm{y^r(\tint)-\yint}{H}^2+ c \norm{y-\Pi^V_{V_r} y}{\mY_T(\tint)}^2\\
      = & \norm{\Pi^H_{V_r} \yint-\yint}{H}^2+c \norm{e_1}{\mY_T(\tint)}^2 \to 0 \ (r\to \infty).
    \end{align*}
    %
    Here we used the fact, that $\Pi^H_{V_r}\yint-\yint \to 0 \ (r\to\infty)$, since $V\subset H$ dense.
    Hence, \eqref{eq:theo:galerkin_conv_help2} implies $e_2\to 0$ in $L^2(V)$ as $r\to\infty$. Next, we show $\norm{\dot e_2}{L^2(V')}=\norm{\dot e_2}{L^2(V_r')}$. We have $y^r\in H^1(V_r)$ and hence $\dot y^r(t)\in V_r\subset V\subset V'\subset V_r'$ for almost all $t\in (\tint,\tint+T)$. Therefore, we conclude the result by applying Riesz's representation theorem
    %
    \begin{equation}\label{eq:hh0}
        \norm{\dot y^r(t)}{V}= \norm{\dot y^r(t)}{V_r}=\norm{\dot y^r(t)}{V_r'}\leq \norm{\dot y^r(t)}{V'}=\norm{\dot y^r(t)}{V}.
    \end{equation}
    %
    Further, we have using $\norm{\Pi^V_{V_r}v }{V}\leq \norm{v}{V}$ for all $v\in V$ and $\Pi_{V_r}^V\circ \Pi_{V_r}^V= \Pi_{V_r}^V $
    %
    \begin{align}\nonumber
       \norm{\dot y(t)\circ  \Pi_{V_r}^V}{V_r'} \leq \norm{\dot y(t)\circ  \Pi_{V_r}^V}{V'}=& \sup_{ v\in V\setminus\{0\}}\tfrac{| \langle \dot y(t), \Pi_{V_r}^Vv\rangle_{V',V}|}{\norm{v}{V}}\leq\sup_{ v\in V\setminus\{0\}}\tfrac{| \langle \dot y(t), \Pi_{V_r}^Vv\rangle_{V',V}|}{\norm{ \Pi_{V_r}^Vv}{V}}\\
        =&\sup_{ v\in V_r\setminus\{0\}}\tfrac{| \langle \dot y(t)\circ  \Pi_{V_r}^V,v\rangle_{V',V}|}{\norm{v}{V}}= \norm{\dot y(t)\circ \Pi_{V_r}^V}{V_r'}.\label{eq:hh1}
    \end{align}
    %
    Hence,
    %
    \begin{align}\nonumber
        \norm{\dot e_2}{L^2(V')}^2 =& \norm{\dot y^r}{L^2(V')}^2-2\langle \dot y^r,  \dot y\circ  \Pi_{V_r}^V \rangle_{L^2(V')}+\norm{\dot y\circ  \Pi_{V_r}^V}{L^2(V')}^2\\
        \stackrel{\eqref{eq:hh0},\eqref{eq:hh1}}{=}&\norm{\dot y^r}{L^2(V_r')}^2-2\langle \dot y^r,  \dot y\circ  \Pi_{V_r}^V \rangle_{L^2(V_r')}+\norm{\dot y\circ \Pi_{V_r}^V}{L^2(V_r')}^2=\norm{\dot e_2}{L^2(V_r')}\label{eq:hh2}
    \end{align}
    %
    Now we can estimate using \eqref{eq:theo:galerkin_conv_help0} as
    %
    \begin{align*}
         \norm{\dot e_2}{L^2(V')}\stackrel{\eqref{eq:hh2}}{=} \norm{\dot e_2}{L^2(V_r')}
         %=\sup_{v\in V_r, \norm{v}{L^2(V)}=1}|\langle \dot e_2,v \rangle_{L^2(V'),L^2(V)}|\\
         &\stackrel{\eqref{eq:theo:galerkin_conv_help0}}{=}\sup_{v\in V_r, \norm{v}{L^2(V)}=1}|\langle \dot e_1+Ae_1-Ae_2,v \rangle_{L^2(V'),L^2(V)}|\\
         & \leq c \norm{\dot e_1}{L^2(V')}+ c\norm{e_1}{L^2(V)}+c\norm{e_2}{L^2(V)}\to 0 \ (r\to \infty).
    \end{align*}
    %
\end{proof}

\end{document}
